\section{Boolean-valued forcing and bare relative consistency}
\label{sec:p06-boolean-forcing-consistency}

A quotient of internally constructed names works even when the initial
model is externally ill-founded. Its elements are equivalence classes of
names, not recursively evaluated sets. Fullness supplies the existential
witness in the quotient truth theorem; the ultrafilter need not be generic.
This route is described in
\cite[Lemmas~1--2 and~10]{HamkinsSeabold2012BooleanV1}.
The proofs below supply the equality, set-bound and axiom arguments needed
here rather than treating that paper's proof sketches as completed proofs.
The separately identified \cite{Bell2005,Jech2003,Kunen2011} presentations
and the historical Scott--Solovay attribution retain their own source duties.

\subsection{Complete Boolean operations and names}

\begin{hbdefinition}{def:p06-bfc-regular-open}{The regular-open completion of set forcing}
For a nonempty poset \(P\), write \(\downarrow p=\{q:q\leq_P p\}\).
An open subset is downward closed. For such \(U\subseteq P\), put
\[
 j(U)=\{p:U\text{ is dense below }p\},\qquad
 n(U)=\{p:\downarrow p\cap U=\varnothing\}.
\]
The regular-open algebra is \(B=\{U\subseteq P:j(U)=U\}\).
Its zero and unit are \(\varnothing,P\), its complement is \(n(U)\),
its meet is intersection and its join is \(j(\bigcup\mathcal A)\)
for a set-indexed family \(\mathcal A\subseteq B\).
The map \(i(p)=j(\downarrow p)\) takes conditions to nonzero elements.
These are ground set constructions, not a poset of external generics.
\end{hbdefinition}

\begin{hbtheorem}{thm:p06-bfc-completion}{proved=annex:p06.bfc.completion}{A nontrivial complete Boolean algebra}
The displayed operations make \(B\) a nontrivial complete Boolean algebra.
The image of \(i\) is dense in \(B\setminus\{0_B\}\); compatibility of
two conditions is equivalent to nonzero meet of their images. Thus ccc of
\(P\) implies ccc of \(B\). All these statements are theorems of ZFC
about ordinary sets and can be interpreted in any model of ZFC.
\end{hbtheorem}

\begin{hbdefinition}{def:p06-bfc-values}{Function-style Boolean names and truth values}
Work first inside ZFC, with a complete Boolean algebra \(B\).
Retain the function-style Boolean names of
\cref{def:p06-boolean-valued-universe}: the domain of a name consists of
earlier names, and its values are in \(B\). Let \(N_B(u)\) be the resulting
first-order name predicate. Write \(\lVert\varphi\rVert_B\) for its
Boolean truth value. The atomic recursion is
\[
\begin{aligned}
 I_B(u,v)&=\bigvee_{t\in\operatorname{dom}(v)}v(t)\wedge_B E_B(u,t),\\
 E_B(u,v)&=\bigwedge_{t\in\operatorname{dom}(u)}
       (\neg_B u(t)\vee_B I_B(t,v))\
       \wedge_B
       \bigwedge_{t\in\operatorname{dom}(v)}
       (\neg_B v(t)\vee_B I_B(t,u)).
\end{aligned}
\]
These assign membership and equality values. For a fixed standard formula,
falsehood has value \(0_B\), negation is Boolean complement, conjunction
is meet, and
\[
 \lVert\exists x\theta(x,\bar u)\rVert_B
 =\bigvee_B\{b\in B:\exists v\ (N_B(v)\land
                             b=\lVert\theta(v,\bar u)\rVert_B)\}.
\]
The set in braces is obtained by Separation on \(B\). No join is indexed
by a primitive proper class. Disjunction, implication and universal
quantification have their fixed classical expansions. This supplies one
defining formula for each fixed formula, not an internal all-code truth
predicate. A missing name coefficient is treated as zero.
\end{hbdefinition}

\begin{hbtheorem}{thm:p06-bfc-values}{proved=annex:p06.bfc.values}{The recursion, equality laws and logical soundness}
The bounded name/truth recursions exist and are compatible on larger bounds.
They are fixed-formula first-order definitions in ZFC. They satisfy
\[
 E_B(u,u)=1_B,\quad E_B(u,v)=E_B(v,u),\quad
 E_B(u,v)\wedge_B E_B(v,w)\leq_B E_B(u,w)
\]
and membership substitution in both arguments. Equality substitution
extends to every fixed formula. Every derivation in the fixed classical
first-order calculus preserves Boolean truth, with the meet of its open
assumption values below the conclusion value. In particular a theorem
from axioms of value one has value one.
\end{hbtheorem}

\begin{hbtheorem}{thm:p06-bfc-fullness}{proved=annex:p06.bfc.fullness}{Mixing and maximal witnesses}
For a set antichain \(A\subseteq B^+\) and ground family \((u_a)_{a\in A}\),
there is a name \(u\) with \(a\leq_B E_B(u,u_a)\) for every \(a\).
For every fixed formula and parameter tuple there is a name \(v\) with
\[
 \lVert\exists x\theta(x,\bar u)\rVert_B
       =\lVert\theta(v,\bar u)\rVert_B.
\]
The antichain, witness family and mixed name are sets. Ground Collection
and Choice are used explicitly; no class-sized choice function is presumed.
\end{hbtheorem}

\begin{hbtheorem}{thm:p06-bfc-axioms}{proved=annex:p06.bfc.axioms}{The ZFC axiom-value theorem scheme}
For every fixed ZFC axiom or schema instance \(\sigma\), ZFC proves
\(\lVert\sigma\rVert_B=1_B\) for every complete Boolean algebra \(B\).
The proofs give pair and union names, separating coefficients, subset-name
bounds, a Collection bound, the minimal-name-rank Foundation argument, and
a ground-indexed function witnessing the Choice construction.
This is a theorem scheme, not a single internal assertion of truth for
all ZFC formula codes.
\end{hbtheorem}

\subsection{An actual quotient of a possibly ill-founded model}

\begin{hbconvention}{def:p06-bfc-ill-founded}{Internal constructions versus external sets}
Let \(\mathcal M=(D,E)\models\mathsf{ZFC}\) be countable. It need not be
well-founded or transitive. Let \(b\in D\) be an object which
\(\mathcal M\) says is a nontrivial complete Boolean algebra.
Its interpreted underlying set
\(B^{\mathcal M}=\{d\in D:\mathcal M\models d\in b\}\)
and operations form an externally countable Boolean algebra.
It need not be externally complete.
The external set of internal names is
\(\mathcal D_b=\{d\in D:\mathcal M\models N_b(d)\}\).
For a fixed formula, its interpreted defining formula assigns an actual
element \(v^{\mathcal M}_\varphi(\bar u)\in B^{\mathcal M}\).
Internal recursion does not authorize an external well-founded evaluation
\(u\mapsto u^G\), and none is used below.
\end{hbconvention}

\begin{hbtheorem}{thm:p06-bfc-quotient}{proved=annex:p06.bfc.quotient}{Quotient truth without external well-foundedness}
For any external ultrafilter \(U\) on \(B^{\mathcal M}\), define
\[
 u\sim_U v\ \Longleftrightarrow\ v^{\mathcal M}_{x=y}(u,v)\in U,
 \qquad
 [u]_U\in_U[v]_U\ \Longleftrightarrow\
                  v^{\mathcal M}_{x\in y}(u,v)\in U.
\]
These define an actual set-sized membership structure
\(\mathcal Q=(\mathcal D_b/{\sim_U},\in_U)\). For every fixed standard
formula,
\[
 \mathcal Q\models\varphi([\bar u]_U)
 \quad\Longleftrightarrow\quad
 v^{\mathcal M}_\varphi(\bar u)\in U.
\]
Consequently \(\mathcal Q\models\mathsf{ZFC}\). Genericity of \(U\),
external completeness of \(B^{\mathcal M}\), and external well-foundedness
of \(\mathcal M\) or \(\mathcal Q\) are not premises.
\end{hbtheorem}

\begin{hbtheorem}{thm:p06-bfc-cohen}{proved=annex:p06.bfc.cohen}{The Boolean non-CH calculation}
Inside ZFC, put \(\kappa=\aleph_2\),
\(P=\operatorname{Fn}(\kappa\times\omega,2,{<}\omega)\), and
\(B=\operatorname{RO}(P)\). Then \(B\) is nontrivial and
\[
 \lVert\neg\mathsf{CH}\rVert_B=1_B.
\]
The proof separately establishes ccc, preservation of the two relevant
ground successor cardinals, and the injection of \(\check\kappa\)
into subsets of \(\check\omega\). One Cohen real is not used as a
substitute for this cardinal-sized family.
\end{hbtheorem}

\begin{hbtheorem}{thm:p06-bfc-consistency}{proved=annex:p06.bfc.consistency}{The bare relative-consistency implication}
In the stated classical metatheory,
\[
 \mathsf{Con}(\mathsf{ZFC})
 \quad\Longrightarrow\quad
 \mathsf{Con}(\mathsf{ZFC}+\neg\mathsf{CH}).
\]
Here \(\mathsf{Con}\) denotes metalinguistic syntactic consistency in
the fixed classical first-order calculus. Completeness supplies an arbitrary
countable model from that premise; it supplies no transitive model.
The interpreted Boolean construction and the actual quotient above supply
the target model. No primitive-recursive PA proof translation or formalized
arithmetical implication is claimed by this metalogical result.
\end{hbtheorem}

\begin{hbexample}{def:p06-bfc-nongeneric}{What an arbitrary ultrafilter preserves}
For \(B=\mathcal P(\omega)\) and a nonprincipal ultrafilter, a mixed name
which is \(\check n\) on the singleton Boolean value \(\{n\}\) has
Boolean value one of being a ground natural number, but its quotient class
need not equal any particular \([\check n]_U\). The join of all singleton
values belongs to \(U\), while none of them does. Fullness, rather than
closure of \(U\) under infinite joins, supplies witnesses in the truth
theorem. The quotient is not claimed externally well-founded or to have
the old model's ordinal height. This is the distinction that the earlier
transitive generic-evaluation presentation cannot silently remove.
\end{hbexample}

\begin{proofboundary}{def:p06-bfc-source-boundary}{Scope of the bridge}
The new proofs do not complete the full original Jech, Kunen or Bell
presentations, or identify a unique Scott--Solovay 1967 manuscript.
The object-theory consistency implication does not establish the consistency
of ZFC without its premise. Exact formalized proof-code transformations,
class forcing and iteration constructions remain separate obligations.
\end{proofboundary}
